\begin{center}
    \LARGE
    \textbf{RELATÓRIO DE MACROENTREGA I}

    \vspace{0.6cm}
    \normalsize

    \begin{minipage}{.90\linewidth}
        \centering
        \textbf{
            Otimização do dispositivo de referência de
monitoramento de movimentos em estruturas
        }
    
        \vspace{.6cm}
        \textbf{PIFF-2609.0073}
        \normalsize
    \end{minipage}
    
    \vspace{1.2cm}
    
    \begin{flushleft}
     \phantom{a.}\textbf{Período: 04/09/2026 a 03/12/2026}
    \end{flushleft}
    
    \begin{center}
    \small
    \begin{tabular}{|c|c|c|}
    \hline
        Item & Produto & Critérios de aceitação \\
    \hline
        \multicolumn{1}{|m{5cm}|}{
            2.1 Mapear lacunas (Gap Analysis) do protótipo
frente às normas IEC e Portaria 115:2022.
        } 
        & \multicolumn{1}{m{3.5cm}|}{
            \centering
            Relatório
        } 
        & \multicolumn{1}{m{6.5cm}|}{
            Relatório contendo uma
planilha de mapeamento de
requisitos normativos
preenchida, identificando
pontos de conformidade e
não-conformidade do
protótipo.

        }\\
    \hline
        \multicolumn{1}{|m{5cm}|}{
            2.2 Modelar em CAD 3D o novo invólucro à
prova de explosão Tipo "d".
        } 
        & \multicolumn{1}{m{3.5cm}|}{
            \centering
            Relatório
        } 
        & \multicolumn{1}{m{6.5cm}|}{
            Disponibilização dos arquivos
de modelagem digital em
formato padrão (ex: .STEP ou
.IGES) acompanhados de desenhos técnicos em PDF
extraídos diretamente do
software.

        }\\
    \hline
        \multicolumn{1}{|m{5cm}|}{
            2.3 Simular estresse mecânico e térmico via
software - métodos clássicos.
        } 
        & \multicolumn{1}{m{3.5cm}|}{
            \centering
            Relatório
        } 
        & \multicolumn{1}{m{6.5cm}|}{
            Apresentação de relatório
descritivo simples contendo as
premissas de cálculo adotadas.

        }\\
    \hline
        \multicolumn{1}{|m{5cm}|}{
            2.4 Adaptar diagramas esquemáticos e o layout
das PCBs com isolamento galvânico.
        } 
        & \multicolumn{1}{m{3.5cm}|}{
            \centering
            Relatório
        } 
        & \multicolumn{1}{m{6.5cm}|}{
            Apresentação da
documentação técnica
atualizada do circuito
eletrônico, acompanhada de
relatório descritivo que
demonstre a implementação
do isolamento galvânico.

        }\\
    \hline
    \end{tabular}
\end{center}

    \vfill
    \centering
    Polo de Inovação\\
    Campos dos Goytacazes\\
    2026

    \newpage

    \vfill
    \textbf{DIRETOR DO PICG}\\
    Vicente de Paulo Santos de Oliveira
    \vfill
    
    \textbf{DIRETOR DE INOVAÇÃO TECNOLÓGICA}\\
    Rogério Atem de Carvalho
    \vfill
    
    \textbf{DIRETORA DE GESTÃO DE PROJETOS}\\
    Juliana Gonçalves Vidigal
    \vfill
    
    \textbf{COORDENADOR DO PROJETO}\\
    Robson da Cunha Santos
    \vfill
    
    \textbf{PESQUISADORES}\\
    Gerson Gomes Cunha\\
    Milena Bissonho Soares\\
    Juliana de Souza Monteiro de Matos\\
    Jean Viana Ribeiro \\
    Klauss Miranda Marchi \\
    Letícia Barcelos Reis \\
    Matheus Ribeiro Vidal
    \vfill
    
    \textbf{ESTUDANTES}\\
    Abraão Carlo Meireles Leite\\
    Gabriel Renato Pereira Campos\\
    Danilo Ricardo Lima de Oliveira\\
    Yuri Florêncio Pereira\\
    Luiza Silva Martins\\
    Davi Gomes Pola\\
    Vitória Cordeiro Mendonça\\
    Lukas Neves de Carvalho
    \vfill

    \textbf{VOLUNTÁRIOS}\\
    Enzo Teixeira Nunes
    \vfill
\end{center}